\documentclass[11pt,a4paper]{article}
\usepackage[T1]{fontenc}
\usepackage{lmodern,microtype,amsmath,amssymb,amsthm,mathtools}
\usepackage[margin=25mm,headheight=14pt]{geometry}
\usepackage{booktabs,array,tabularx,enumitem,xcolor,fancyhdr,listings,hyperref}
\definecolor{ink}{HTML}{17324D}
\definecolor{accent}{HTML}{176B78}
\definecolor{codebg}{HTML}{F5F7F9}
\hypersetup{colorlinks=true,linkcolor=ink,citecolor=accent,urlcolor=accent,
pdftitle={Integral palette lifting in arbitrary degree: a constructive proof of the alpha+n conjecture},
pdfsubject={Hermite cancellation, unimodular nested palettes, and integral chromatic factors}}
\newtheorem{theorem}{Theorem}[section]
\newtheorem{lemma}[theorem]{Lemma}
\newtheorem{proposition}[theorem]{Proposition}
\newtheorem{corollary}[theorem]{Corollary}
\theoremstyle{definition}
\newtheorem{definition}[theorem]{Definition}
\newtheorem{remark}[theorem]{Remark}
\newcommand{\Z}{\mathbb Z}
\newcommand{\Q}{\mathbb Q}
\newcommand{\R}{\mathbb R}
\newcommand{\fall}[2]{(#1)_{#2}}
\newcommand{\He}{\operatorname{He}}
\newcommand{\htf}{\mathcal H}
\newcommand{\calP}{\mathcal P}
\newcommand{\Split}{\operatorname{Split}_{\Q}}
\newcommand{\wt}{\gamma}
\pagestyle{fancy}\fancyhf{}
\fancyhead[L]{\small\color{ink}Integral palette lifting in arbitrary degree}
\fancyhead[R]{\small 28 November 2026}
\fancyfoot[C]{\small\thepage}
\setlength{\emergencystretch}{3em}
\setlist{nosep,leftmargin=*}
\raggedbottom
\lstset{language=Python,basicstyle=\ttfamily\scriptsize,
keywordstyle=\color{ink}\bfseries,commentstyle=\color{accent},
backgroundcolor=\color{codebg},frame=single,rulecolor=\color{black!12},
breaklines=true,columns=fullflexible,keepspaces=true,
showstringspaces=false,tabsize=4,aboveskip=10pt,belowskip=10pt}
\title{\color{ink}\bfseries Integral palette lifting in arbitrary degree\\[3pt]
\large A constructive proof of the \texorpdfstring{$\alpha+n$}{alpha+n} conjecture}
\author{\normalsize Research manuscript for independent verification}
\date{28 November 2026}

\begin{document}
\maketitle
\begin{abstract}
For every monic polynomial $f\in\Z[x]$ of degree $d\ge1$, we construct
positive integers $N,n$ and a connected graph $G$, consisting of an
$N$-clique and a $d$-clique with specified cross-edges, such that
$P_G(q)=\fall qN f(q-n)$ and $\chi(G)=N$.
The construction combines a strictly positive Boolean intersection profile
with three uniform mechanisms. Hermite roots cancel the leading matching
polynomial to $X^d$. Restricting all corrections to intersections containing
one distinguished row makes the coefficient map exactly affine. Finally,
nested palettes supply a coefficient matrix of determinant one modulo
$(d-1)!$, permitting a descending sequence of integral corrections.
All corrections have size $O_d(\htf t)$, whereas every palette atom has a
positive leading term of order $t^2$, where
$\htf=1+\sum_{j=1}^d|a_j|$ for $f=x^d+\sum_{j=1}^d a_jx^{d-j}$.
This proves termination of an exact algorithm and gives
$|V(G)|\le C_d\htf^2$ for a constant depending only on $d$.
The factor identity implies Cameron--Morgan's original $\alpha+n$
statement in every degree and Gilmore's prescribed-degree splitting-field
conjecture. We include a rational alternative to the Hermite profile,
a complete sextic certificate, and a reference implementation with an
independent matching check. The proof is presented for external scrutiny;
the literature search and its limits are recorded separately from the
mathematical argument.
\end{abstract}

\section{The original problem and the result}

For a finite simple graph $G$, let $P_G(q)$ denote its chromatic polynomial.
A \emph{chromatic root} is a complex zero of such a polynomial.
Cameron and Morgan's Conjecture~4 asks whether every algebraic integer
$\alpha$ has an integer translate $\alpha+n$, with $n$ natural, that is a
chromatic root \cite{CM}. There is no restriction on the degree or on the
number of real conjugates of $\alpha$.

Write $\fall qN=q(q-1)\cdots(q-N+1)$, with $\fall q0=1$.
A \emph{$(d,N)$-biclique} in the terminology of the cited literature is
a graph partitioned into cliques of orders $d$ and $N$, with arbitrary
cross-edges. This is the complement of a bipartite graph; it need not be a
complete bipartite graph.

\begin{theorem}[Exact realization in every degree]\label{thm:main}
For every integer $d\ge1$ and every monic polynomial
\[
 f(x)=x^d+a_1x^{d-1}+\cdots+a_d\in\Z[x],
\]
there exist positive integers $N,n$ and a connected finite simple
$(d,N)$-biclique $G$ such that
\begin{equation}\label{eq:main}
 P_G(q)=\fall qN f(q-n),\qquad \chi(G)=N.
\end{equation}
An exact algorithm producing $N,n$ and a finite description of $G$
terminates for every input. For every fixed $d$, there is a constant
$C_d>0$ such that the output satisfies
\begin{equation}\label{eq:size}
 |V(G)|=N+d\le C_d\htf^2,
 \qquad \htf=1+\sum_{j=1}^d|a_j|.
\end{equation}
No irreducibility, signature, coefficient-sign, or Galois-group assumption
is imposed.
\end{theorem}

\begin{corollary}[The original $\alpha+n$ statement]\label{cor:alpha}
For every algebraic integer $\alpha$, there is a positive integer $n$
and a connected finite simple graph $G$ with $P_G(\alpha+n)=0$.
\end{corollary}
\begin{proof}
The monic minimal polynomial $f\in\Z[x]$ of $\alpha$ satisfies the
hypotheses of Theorem~\ref{thm:main}. Substitute $q=\alpha+n$ in
\eqref{eq:main}.
\end{proof}

\subsection{Prior work and the auxiliary problem}

The quadratic case was established by the original study group, and
Bohn proved the cubic case \cite{CM,Bohn}. Gilmore's thesis contains
quartic and quintic computational evidence and formulates a
prescribed-degree splitting-field question \cite[Chapter 7]{Gilmore}.
The official abstract of Delbourgo--Morgan's 2019 paper describes an
algorithm in degrees at most four, conditional on existence of an
appropriate biclique \cite{DM19}. Their 2022 paper calculates the fifth
chromatic coefficient and discusses the earlier evidence \cite{DM22}.
The source-access limitations are stated in Section~\ref{sec:status}.

The auxiliary problem here is to prescribe the entire palette polynomial
while every one of the $2^d-1$ nonempty Boolean atoms remains a positive
integer. Its solution has two logically separate parts: exact integral
coefficient completion, and a uniform size estimate ensuring that this
completion stays inside the positive region. The partition formula alone
does not provide either conclusion.

The supplied corpus manuscript on roots and coefficient peaks of graph
polynomials uses descending coefficient corrections with a positive
margin \cite[Sections 2--3]{VMRC}. That organizing idea motivates the
present construction. The proof below establishes all necessary lemmas
directly and assumes no theorem asserted in that corpus. The earlier
degree-four and degree-five working constructions are likewise not
prerequisites for this argument.

\section{Palette atoms and an affine coefficient map}\label{sec:palettes}

Fix $d\ge2$, and put $[d]=\{1,\ldots,d\}$. For every nonempty
$J\subseteq[d]$, take $m_J$ vertices of type $J$ in a clique $C$.
Set $N=\sum_{J\ne\varnothing}m_J$. Add a disjoint clique
$\{v_1,\ldots,v_d\}$, and make a vertex of type $J$ nonadjacent to
$v_i$ exactly when $i\in J$. Thus $A_i\subseteq C$, the set of
nonneighbors of $v_i$, has intersection moments
\begin{equation}\label{eq:moments}
 H_B=\left|\bigcap_{i\in B}A_i\right|
     =\sum_{J\supseteq B}m_J\qquad(\varnothing\ne B\subseteq[d]).
\end{equation}
Conversely, integer moments determine the atoms by Boolean inversion:
\begin{equation}\label{eq:inversion}
 m_J=\sum_{B\supseteq J}(-1)^{|B|-|J|}H_B.
\end{equation}
If all these integers are nonnegative, the type construction realizes
every moment exactly. There are no additional set-realizability
conditions. We will make every $m_J$ strictly positive.

Choose an integer $L$ and write
\begin{equation}\label{eq:parameters}
 H_{\{i\}}=L+z_i,\qquad H_B=L-D_B\quad(|B|\ge2),
 \qquad c_{\{i\}}=z_i,\quad c_B=-D_B\quad(|B|\ge2).
\end{equation}
For any $R\subseteq[d]$, define the formal palette polynomial
\begin{equation}\label{eq:palette}
 F_R(x)=\sum_{\pi\in\calP(R)}\prod_{B\in\pi}
       \mu_B(x+c_B),\qquad
 \mu_B=(-1)^{|B|-1}(|B|-1)!,\qquad F_\varnothing=1.
\end{equation}
Here $\calP(R)$ is the set of partitions of $R$. The formula is meaningful
for arbitrary integer parameters, including intermediate choices that
do not yet give nonnegative atoms.

\begin{lemma}[The graph identity]\label{lem:graph}
Whenever the atoms in \eqref{eq:inversion} are nonnegative integers,
the graph just described satisfies
\begin{equation}\label{eq:graph}
 P_G(q)=\fall qN F_{[d]}(q-N+L).
\end{equation}
Moreover, $F_R$ is monic of degree $|R|$, and
\begin{equation}\label{eq:trace}
 [x^{d-1}]F_{[d]}=\sum_{i=1}^d z_i-\binom d2.
\end{equation}
\end{lemma}
\begin{proof}
Color $C$ first, using $N$ distinct colors. For $q\ge N$, each $v_i$
can use the $q-N$ unused colors and the old colors corresponding to
$A_i$. The intersection of the available color lists for a nonempty
block $B$ therefore has size $q-N+H_B=x+c_B$, where $x=q-N+L$.
The number of assignments giving pairwise different colors to the
vertices indexed by $R$ is \eqref{eq:palette}.

For completeness, this last assertion can be verified by summing over
permutations of $R$. Give a permutation the weight
$(-1)^{|R|-\text{number of cycles}}$, and require an assignment to be
constant on each cycle. For a fixed assignment the sum factors over its
color fibers. A fiber of size $s$ contributes the sum of the signs of
all permutations in $S_s$, which is $1$ for $s=1$ and $0$ for $s\ge2$.
Thus only injective assignments survive. Grouping permutations by their
cycle partition gives the factor $(-1)^{|B|-1}(|B|-1)!$ for each block.
Multiplication by $\fall qN$ proves \eqref{eq:graph} for all integers
$q\ge N$, hence as a polynomial identity.

Only the singleton partition contributes degree $|R|$. For
$R=[d]$, the degree $d-1$ term comes from the singleton partition
and from the $\binom d2$ partitions with one pair and all other blocks
singletons. The latter each contribute $-1$, giving \eqref{eq:trace}.
\end{proof}

\begin{lemma}[One distinguished row]\label{lem:affine}
Keep $z_1,\ldots,z_d$ and all $D_B$ with $d\notin B$ fixed.
The polynomial $F_{[d]}$ is jointly affine in the variables
$D_B$ with $d\in B$ and $|B|\ge2$. For such a block of size $k$,
\begin{equation}\label{eq:derivative}
 \frac{\partial F_{[d]}}{\partial D_B}
   =\eta_k F_{[d]\setminus B},\qquad
 \eta_k=(-1)^k(k-1)!.
\end{equation}
All polynomials on the right depend only on the fixed outer rows
$\{1,\ldots,d-1\}$.
\end{lemma}
\begin{proof}
Every partition has exactly one block containing $d$.
Consequently it contains at most one of the variables under discussion.
For a specified block $B$, the remaining blocks form an arbitrary
partition of $[d]\setminus B$, and differentiation of
$\mu_B(x-D_B)$ contributes $-\mu_B=\eta_k$.
\end{proof}

In particular, simultaneous changes of these variables have exactly the
sum of their individual effects. There is no error term involving a
product of changes. The estimates later in the proof rely on this exact
property.

\section{Hermite cancellation and uniform rank}\label{sec:hermite}

On the finite-dimensional space of polynomials, let
\[
 T=\exp\!\left(\tfrac12\partial_X^2\right),\qquad
 \He_d(X)=T^{-1}X^d
 =\sum_{j=0}^{\lfloor d/2\rfloor}
 \frac{(-1)^j d!}{2^j j!(d-2j)!}X^{d-2j}.
\]
The exponential is a finite sum on every polynomial, and its inverse is
$\exp(-\partial_X^2/2)$.

\begin{lemma}[The Hermite slopes]\label{lem:hermite}
The polynomial $\He_d$ has $d$ distinct real roots
$\lambda_1<\cdots<\lambda_d$. Their sum is zero, and
\begin{equation}\label{eq:hermite-cancel}
 \prod_{i=1}^d(X+\lambda_i)=\He_d(X),\qquad
 T\prod_{i=1}^d(X+\lambda_i)=X^d.
\end{equation}
\end{lemma}
\begin{proof}
The displayed coefficient formula gives the recurrence
$\He_{r+1}=X\He_r-r\He_{r-1}$, with $\He_0=1$, $\He_1=X$.
It follows by expanding a tridiagonal determinant that $\He_d$ is the
characteristic polynomial of the real symmetric $d\times d$ matrix with
zero diagonal and off-diagonal entries $\sqrt{1},\sqrt{2},\ldots,
\sqrt{d-1}$. Its eigenvalues are real. An eigenvector is determined by
its first coordinate through the successive row equations, and a zero
first coordinate forces the whole vector to be zero. Every eigenspace
is therefore one-dimensional. Symmetric diagonalizability implies that
all eigenvalues are simple.
The coefficient of $X^{d-1}$ is zero, so their sum is zero. The parity
$\He_d(-X)=(-1)^d\He_d(X)$ makes the root set invariant under negation,
which gives the first identity in \eqref{eq:hermite-cancel}; the second
follows from the definition of $T$.
\end{proof}

For a finite index set $R$, put
\begin{equation}\label{eq:leading-Q}
 Q_R(X)=T\prod_{i\in R}(X+\lambda_i).
\end{equation}
These polynomials admit a useful counting expansion:
\begin{equation}\label{eq:heat-matching}
 Q_R(X)=\sum_{\mathcal M\text{ a matching on }R}
             \prod_{i\text{ unmatched by }\mathcal M}(X+\lambda_i).
\end{equation}
Indeed, applying $\partial_X^{2j}/(2^j j!)$ to the product first selects
$2j$ distinct factors and then gives the number
$(2j)!/(2^j j!)$ of pairings of those factors.

\begin{lemma}[Every leading correction matrix is invertible]
\label{lem:rank}
For any $I\subseteq[d-1]$ of size $r+1$, the $r+1$ polynomials
\[
 Q_{I\setminus\{i\}}(X),\qquad i\in I,
\]
are a basis of the polynomials of degree at most $r$.
\end{lemma}
\begin{proof}
Before applying $T$, these are the products
$\prod_{j\in I\setminus\{i\}}(X+\lambda_j)$.
Evaluation at $X=-\lambda_i$ kills every product except the $i$th,
whose value is nonzero because the slopes are distinct. The products
are independent. Applying the invertible linear map $T$ preserves
independence.
\end{proof}

\section{An integral seed inside the positive region}\label{sec:seed}

Fix $d\ge2$ and the input polynomial $f$. Throughout the remaining
construction let
\begin{equation}\label{eq:scale}
 m=d-1,\qquad M=m!,\qquad U=2^dM,\qquad t\in U\Z_{>0},
 \qquad L=2t^2,\qquad \wt=2^{2-d}.
\end{equation}
Choose integers $z_i$ for $1\le i\le m$ such that
\begin{equation}\label{eq:singleton-round}
 z_i\equiv i-1\pmod M,\qquad
 |z_i-t\lambda_i|\le2M.
\end{equation}
This is an effective finite operation. For example, isolate each simple
root in a rational interval of width at most $M/(4t)$, use its midpoint
$b_i$, and set
\[
 z_i=(i-1)+M\left\lfloor\frac{tb_i-(i-1)}{M}\right\rfloor.
\]
The midpoint contributes an error of at most $M/8$, and rounding
contributes less than $M$, so \eqref{eq:singleton-round} follows.
Rational isolation of roots of an integer polynomial can be carried out
by Sturm division and rational bisection; no numerical oracle is needed.
Now set
\begin{equation}\label{eq:center-trace}
 z_d=a_1+\binom d2-\sum_{i=1}^m z_i.
\end{equation}
This fixes the coefficient of $x^{d-1}$ exactly.

For a nonempty set $B$ of size $k\ge2$, initialize its deficit by
\begin{equation}\label{eq:initial-D}
 D_B^{(0)}=(2-2^{2-k})t^2-\varepsilon_B,\qquad
 \varepsilon_B=
 \begin{cases}
  \min B-1,&d\notin B,\\
  0,&d\in B.
 \end{cases}
\end{equation}
Every quadratic term here is an integer multiple of $M$ by
\eqref{eq:scale}. For all nonempty outer sets $B\subseteq[m]$, the
parameters in \eqref{eq:parameters} consequently satisfy
\begin{equation}\label{eq:nested-congruence}
 c_B\equiv \min B-1\pmod M.
\end{equation}

The quadratic part of these moments is especially simple. Giving every
nonempty atom the weight $\wt t^2$ gives singleton moment $2t^2$ and
moment $2^{2-|B|}t^2$ on every nonempty set $B$. Thus the initialized
moments differ from this strictly positive uniform profile only in terms
of order at most $t$ for fixed $f$.

\begin{lemma}[Uniform scaled limits]\label{lem:limits}
Use the seed just specified, and keep the outer data fixed thereafter.
For each outer set $R$ of size $r$,
\begin{equation}\label{eq:outer-limit}
 t^{-r}F_R(tX)=Q_R(X)+O_d(t^{-1})
\end{equation}
coefficient by coefficient, uniformly over all choices permitted in
\eqref{eq:singleton-round}. For the initial full polynomial,
\begin{equation}\label{eq:full-limit}
 t^{-d}F_{[d]}^{(0)}(tX)=X^d+O_d(\htf/t).
\end{equation}
After its leading and next-to-leading coefficients are matched, the
initial residual $E^{(0)}=f-F_{[d]}^{(0)}$ has degree at most $d-2$ and
\begin{equation}\label{eq:initial-error}
 |[x^j]E^{(0)}|\le A_d\htf\,t^{d-1-j}
 \qquad(0\le j\le d-2)
\end{equation}
for some constant $A_d$ independent of $f$ and $t\ge U$.
\end{lemma}
\begin{proof}
In a singleton block, $tX+z_i=t(X+\lambda_i)+O_d(1)$ for an
outer row. A pair block contributes
$-(tX-D_B^{(0)})=t^2+O_d(t)$.
A block of size $k\ge3$ contributes a polynomial of size $O_d(t^2)$
coefficientwise in $X$, which has strictly smaller order than $t^k$.
In the sum over partitions of $R$, only singleton-pair partitions have
order $t^r$. Their sum is \eqref{eq:heat-matching}; every other partition
has order at most $t^{r-1}$. All errors are uniform because the rounding
errors lie in fixed bounded intervals.

For the full set, \eqref{eq:center-trace} and
$\sum_i\lambda_i=0$ give
$z_d=t\lambda_d+O_d(\htf)$. There is only one block containing row $d$
in each partition, so the dependence on this larger error is linear.
The same expansion therefore has error $O_d(\htf t^{d-1})$.
Equation~\eqref{eq:hermite-cancel} gives the leading polynomial $X^d$.
Both leading coefficients match $f$ by monicity and
\eqref{eq:center-trace}. Comparing coefficients of $X^j$ in
$f(tX)-F_{[d]}^{(0)}(tX)$ proves \eqref{eq:initial-error}.
\end{proof}

\section{Nested palettes remove the lattice obstruction}\label{sec:lattice}

For $s\ge1$, consider the formal outer model $c_B=\min B-1$ on
$[s]$. At every sufficiently large integer $x$, realize it by nested
color lists of sizes $x,x+1,\ldots,x+s-1$.

\begin{lemma}[A nested basis of determinant one]\label{lem:nested}
For $1\le i\le s$ in this model,
\begin{equation}\label{eq:nested-basis}
 F_{[s]\setminus\{i\}}(x)=x^{i-1}(x+1)^{s-i}.
\end{equation}
These $s$ polynomials have an integer coefficient matrix of determinant
$1$ in the ascending monomial basis $1,x,\ldots,x^{s-1}$.
For the actual seed, the same identities hold modulo $M$ for every
prefix $[s]\subseteq[m]$.
\end{lemma}
\begin{proof}
Color the remaining rows in increasing order. Before the omitted row,
row $j$ has a list of size $x+j-1$ and $j-1$ previously used colors,
all lying in its list, leaving $x$ choices. After the omitted row it
has only $j-2$ previously used colors, leaving $x+1$ choices. This gives
\eqref{eq:nested-basis} on infinitely many integer values of $x$, hence
identically. The $i$th column starts with $i-1$ zeros and has coefficient
$1$ at $x^{i-1}$, so the coefficient matrix is triangular with diagonal
$1$. Finally, \eqref{eq:palette} has integer coefficients in the
parameters $c_B$; reducing \eqref{eq:nested-congruence} modulo $M$
therefore gives the congruence assertion.
\end{proof}

Let the current residual be $E=f-F_{[d]}$. Initially $\deg E\le d-2$.
Changing the star deficits $D_{\{i,d\}}$, for $1\le i\le m$, changes
$F_{[d]}$ by the columns $F_{[m]\setminus\{i\}}$, since $\eta_2=1$.
Lemma~\ref{lem:nested} makes their coefficient matrix invertible
modulo $M$. Choose residues
\begin{equation}\label{eq:rho}
 \rho_i\in\{0,1,\ldots,M-1\}
\end{equation}
so that replacing $D_{\{i,d\}}$ by $D_{\{i,d\}}+\rho_i$ makes every
coefficient of $E$ divisible by $M$. For $M=1$ take every $\rho_i=0$.
These bounded changes do not affect the leading or next-to-leading
coefficients.

We now use blocks of sizes $k=2,3,\ldots,d$, in that order.
At stage $k$, put
\begin{equation}\label{eq:stage-blocks}
 r=d-k,\qquad I=[r+1],\qquad
 R_i=I\setminus\{i\},\qquad B_i=[d]\setminus R_i
 \quad(1\le i\le r+1).
\end{equation}
Each $B_i$ contains $d$ and has size $k$. Its derivative column is
$\eta_k F_{R_i}$, a polynomial of degree $r$ with leading coefficient
$\eta_k$.

\begin{lemma}[Integral completion with controlled rounding]
\label{lem:completion}
Suppose at the start of stage $k$ that $\deg E\le r$ and
$E\in M\Z[x]$. There are integral changes of the deficits indexed by
\eqref{eq:stage-blocks} that remove the coefficient of $x^r$ exactly
and leave the entire residual in $M\Z[x]$.
They differ by a bounded amount depending only on $d$ from the unique
rational changes that would remove all remaining coefficients at once.
After stage $d$ the residual is zero.
\end{lemma}
\begin{proof}
The matrix with columns $F_{R_i}$ has determinant $1$ modulo $M$ by
Lemma~\ref{lem:nested}. For $d\ge3$ this determinant is nonzero, because
$M\ge2$. For $d=2$ the only matrix is $[1]$. Thus the matrix with columns
$\eta_k F_{R_i}$ is invertible over $\Q$ at every admissible scale.
Let $\delta_1^*,\ldots,\delta_{r+1}^*$ solve
\begin{equation}\label{eq:exact-stage}
 E(x)=\sum_{i=1}^{r+1}\delta_i^*\eta_kF_{R_i}(x).
\end{equation}
Put
\begin{equation}\label{eq:lattice-step}
 h_k=\frac{M}{(k-1)!}.
\end{equation}
For the first $r$ coordinates choose
\begin{equation}\label{eq:round-stage}
 \widehat\delta_i=h_k\left\lfloor\delta_i^*/h_k\right\rfloor,
 \qquad
 \widehat\delta_{r+1}=
 \frac{[x^r]E}{\eta_k}-\sum_{i=1}^r\widehat\delta_i.
\end{equation}
Since $[x^r]E$ is divisible by $M$, its quotient by $\eta_k$ lies in
$h_k\Z$. Therefore every $\widehat\delta_i$ belongs to $h_k\Z$.
All column leading coefficients equal $\eta_k$, so
\eqref{eq:round-stage} removes the degree $r$ term exactly.
Also $\eta_kh_k=\pm M$ and $F_{R_i}\in\Z[x]$, which preserves
divisibility of every residual coefficient by $M$.

For $i\le r$, the error
$e_i=\widehat\delta_i-\delta_i^*$ belongs to $(-h_k,0]$.
Taking leading coefficients in \eqref{eq:exact-stage} shows
$e_{r+1}=-\sum_{i=1}^re_i$. Hence
\begin{equation}\label{eq:round-error}
 |e_i|\le r h_k\quad(1\le i\le r+1),\qquad \sum_i e_i=0
\end{equation}
when $r\ge1$. At the final stage $r=0$, the single derivative is the
constant $\eta_d=\pm M$, and the remaining constant is removed exactly
with an integer change; the rounding error is zero.
All changes are additive by Lemma~\ref{lem:affine}, so the stage
calculation is exact even after the preceding stages.
\end{proof}

\section{The corrections stay inside the positive region}\label{sec:bounds}

We now establish the quantitative assertion needed to turn the formal
integer solution into a graph. In all estimates below, $d$ is fixed,
and the constants are independent of the input coefficients and of
the permitted root-rounding choices.

\begin{lemma}[Uniform bound for every correction]\label{lem:bound}
There is a constant $K_d>0$ such that, for every input $f$ and
every admissible $t\ge U$, the residue repair and all stages of
Lemma~\ref{lem:completion} change each initialized deficit by at most
$K_d\htf t$ in absolute value. The outer deficits remain fixed.
\end{lemma}
\begin{proof}
At stage $k$, with $r=d-k$, define the scaled matrix $A_k(t)$ by
\begin{equation}\label{eq:scaled-matrix}
 (A_k(t))_{j,i}=t^{j-r}[x^j](\eta_k F_{R_i}(x)),
 \qquad 0\le j\le r,\quad 1\le i\le r+1.
\end{equation}
By Lemma~\ref{lem:limits}, this converges uniformly to the coefficient
matrix with columns $\eta_kQ_{R_i}$. The latter is invertible by
Lemma~\ref{lem:rank}. Thus, beyond a degree-dependent threshold, all
inverses $A_k(t)^{-1}$ are bounded in any fixed matrix norm.
There are only $d-1$ stages, so one threshold and one bound suffice.
Below that threshold there are only finitely many admissible scales
and finitely many permitted integer outer data. Every corresponding
matrix is nonsingular by Lemma~\ref{lem:completion}. Increasing the
degree-dependent bound over these finitely many matrices therefore
bounds all the scaled inverses for every admissible $t\ge U$.
The coefficients $t^{j-r}[x^j]F_{R_i}$ are likewise uniformly bounded.

For the first stage, \eqref{eq:initial-error} gives
$[x^j]E=O_d(\htf t^{r+1-j})$.
The bounded residue repair contributes only $O_d(t^{r-j})$ and
preserves this bound. Multiplying the right side of
\eqref{eq:exact-stage} by the row scalings in
\eqref{eq:scaled-matrix} gives a vector $O_d(\htf t)$.
The inverse bound therefore gives
$\delta_i^*=O_d(\htf t)$ in this first stage.

After any stage of residual degree $r\ge1$, the new residual is
\[
 -\sum_i e_i\eta_kF_{R_i}(x).
\]
By \eqref{eq:round-error} its coefficient at $x^j$ is
$O_d(t^{r-j})$, and its coefficient at $x^r$ is exactly zero.
At the next stage the degree is at most $r-1$; scaling the rows for
that degree makes this new right side $O_d(t)$.
Every later exact correction is consequently $O_d(t)$, and its rounded
version has the same bound. The final constant stage satisfies the
same assertion and removes its residual without rounding error.

The initial residue changes are bounded by $M$, and there are finitely
many stages and coordinates depending only on $d$. Summing the changes
gives $O_d(\htf t)$ for each affected deficit. All blocks changed
contain $d$, so all outer data used in the matrices stay fixed.
\end{proof}

\begin{lemma}[Positive atoms and a positive shift]\label{lem:positive}
For the final moments after integral completion,
\begin{align}
 m_J&=\wt t^2+O_d(\htf t)\qquad(\varnothing\ne J\subseteq[d]),
     \label{eq:atoms-asymp}\\
 N&=(4-\wt)t^2+O_d(\htf t),\qquad
 n=N-L=(2-\wt)t^2+O_d(\htf t).\label{eq:N-asymp}
\end{align}
There is a constant $B_d>0$ such that all final atoms and $n$ are
strictly positive whenever $t\ge B_d\htf$ is admissible.
\end{lemma}
\begin{proof}
Before the corrections, the quadratic part of every moment comes from
the uniform atom weights $\wt t^2$. The singleton perturbations are
$O_d(t+\htf)$ and the initial outer offsets are $O_d(1)$.
Lemma~\ref{lem:bound} changes the other moments by
$O_d(\htf t)$. Boolean inversion is a fixed finite linear transform
for this degree, proving \eqref{eq:atoms-asymp} uniformly for every
nonempty $J$. Summing its $2^d-1$ terms gives
$(2^d-1)\wt=4-\wt$, proving \eqref{eq:N-asymp}.

For $d\ge2$, both $\wt=2^{2-d}$ and $2-\wt$ are positive.
Choose $B_d\ge U$ large enough that the absolute error bounds in
\eqref{eq:atoms-asymp} and \eqref{eq:N-asymp} are smaller than half
their respective positive leading terms whenever $t\ge B_d\htf$.
This is possible because every such error is bounded by a constant
depending only on $d$ times $\htf t$. It proves simultaneous positivity
of all atoms and of $n$ for all sufficiently large admissible scales.
\end{proof}

\begin{proof}[Proof of Theorem~\ref{thm:main}]
Suppose first that $d\ge2$. Starting with $t=U$, successively try
$t=U,2U,4U,\ldots$. At each scale carry out the finite root isolation,
the residue repair, and all the exact rational systems and lattice
roundings above. Accept the first scale at which all computed atoms
and $n=N-L$ are positive. Lemma~\ref{lem:positive} proves that this
unbounded doubling search terminates for every input. At the accepted
scale, Lemma~\ref{lem:completion} gives $F_{[d]}=f$ exactly, and the
positive integer atoms give an actual graph. Lemma~\ref{lem:graph}
then gives
\[
 P_G(q)=\fall qN f(q-N+L)=\fall qN f(q-n).
\]

The graph contains the $N$-clique $C$. For each $i$, choose a vertex of
the positive singleton type $\{i\}$ and reuse its color on $v_i$.
These are distinct colors and each is allowed on its corresponding
row, so $\chi(G)=N$. A vertex of type $\{1\}$ is adjacent to $v_2$;
there is therefore an edge joining the two cliques. Both cliques are
connected, so the graph is connected.

The first accepted scale is at most
$2\max(U,B_d\htf)$, after enlarging $B_d$ if necessary.
Equation~\eqref{eq:N-asymp} now gives
$N+d=O_d(\htf^2)$, which is \eqref{eq:size}.

For $d=1$, write $f(x)=x+a_1$ and take
\[
 N=\max(2,2-a_1),\qquad n=N+a_1-1\ge1,\qquad
 G=K_{N+1}-e.
\]
Partition the graph into an $N$-clique and a single vertex that is
nonadjacent to exactly one vertex of the clique. Direct counting gives
\[
 P_G(q)=\fall qN(q-N+1)=\fall qN f(q-n).
\]
This graph is connected, has chromatic number $N$, and has
$O(\htf)$ vertices. This completes all cases.
\end{proof}

\begin{remark}[Quantifiers and effectiveness]\label{rem:quantifiers}
For each $d\ge2$, one degree-dependent constant $B_d$ works for every
input $f$ and for every admissible scale $t\ge B_d\htf$.
Thus the positivity argument proves existence
for every input and termination of the stated algorithm, rather than
giving a sufficient condition whose attainability is left open.
The constants in the size bound are degree dependent; no bound uniform
in unbounded $d$ is asserted. They can be made effective from the finite
partition formulas and certified bounds for the distinct algebraic
slopes and the inverses of their limiting matrices. Their numerical
optimization is not needed by the algorithm or the existence proof.
\end{remark}

\section{Prescribed-degree splitting fields}\label{sec:fields}

Gilmore's Conjecture~7.1 asks whether a field that is the splitting field
of a degree-$j$ polynomial over $\Q$ is the splitting field of the
chromatic polynomial of a $(j,k)$-biclique for some $k$
\cite[p.~63]{Gilmore}. The parameter $j$ is the degree of the given
polynomial, not the degree of its splitting field.

\begin{corollary}[Exact splitting-field realization]\label{cor:field}
Let $g\in\Q[x]$ have degree $d\ge1$. There is a connected
$(d,N)$-biclique $G$ such that
\[
 \Split(P_G)=\Split(g).
\]
Consequently the prescribed-degree statement in Gilmore's
Conjecture~7.1 follows from Theorem~\ref{thm:main}.
\end{corollary}
\begin{proof}
Divide $g$ by its leading coefficient, obtaining
$g_0(x)=x^d+b_1x^{d-1}+\cdots+b_d$.
Choose a positive common denominator $s$ of the rational $b_i$.
Then
\[
 f(x)=s^d g_0(x/s)
     =x^d+s b_1x^{d-1}+\cdots+s^d b_d\in\Z[x]
\]
is monic of degree $d$ and has exactly the same splitting field as $g$:
its roots are $s$ times the roots of $g$.
Theorem~\ref{thm:main} gives
$P_G(q)=\fall qN f(q-n)$ with a small clique of exactly $d$ vertices.
Translation by a rational integer and adjoining the rational roots of
$\fall qN$ do not change the splitting field. Hence the asserted
equality holds.
\end{proof}

This realizes splitting fields already arising from polynomials over
$\Q$. It does not assert that every finite abstract group occurs as a
Galois group over $\Q$.

\section{An exact certificate and independent verification}\label{sec:certificate}

The reference implementation in Appendix~\ref{app:code} accepts the
coefficient list $[a_1,\ldots,a_d]$. It uses rational isolating intervals
only to choose the auxiliary Hermite slopes; the original target roots
are never approximated. All subsequent arithmetic is integral or
rational, and each acceptance test is exact.

\subsection{A complete sextic certificate}

For $f(x)=x^6-2$, the implementation returns
\begin{align*}
 t&=15360,&L&=471859200,&N&=951531708,\\
 n&=479672508,&|V(G)|&=951531714,&\min_J m_J&=3894438,
\end{align*}
and
\[
 (z_1,\ldots,z_6)=(-51120,-29039,-9478,9363,28924,51365).
\]
Initialize every deficit by \eqref{eq:initial-D} and then add the
following corrections; all unlisted corrections are zero.
The moments \eqref{eq:parameters} and inversion
\eqref{eq:inversion} specify every vertex type and its multiplicity,
so the certificate is a complete compressed description of the graph.

\begin{center}
\small
\begin{tabular}{@{}lr@{\qquad}lr@{}}
\toprule
$B$ & $D_B-D_B^{(0)}$ & $B$ & $D_B-D_B^{(0)}$\\
\midrule
$\{1,6\}$ & $6\,662\,758$ & $\{4,5,6\}$ & $-193\,500$ \\
$\{2,6\}$ & $1\,849\,209$ & $\{1,4,5,6\}$ & $-551\,280$ \\
$\{3,6\}$ & $339\,460$ & $\{2,4,5,6\}$ & $-369\,520$ \\
$\{4,6\}$ & $4\,592\,431$ & $\{3,4,5,6\}$ & $-127\,780$ \\
$\{5,6\}$ & $2\,100\,692$ & $\{1,3,4,5,6\}$ & $-187\,975$ \\
$\{1,5,6\}$ & $-4\,560\,780$ & $\{2,3,4,5,6\}$ & $-10\,645$ \\
$\{2,5,6\}$ & $-1\,656\,900$ & $\{1,2,3,4,5,6\}$ & $-9\,071$ \\
$\{3,5,6\}$ & $-1\,462\,200$ &  &  \\
\bottomrule
\end{tabular}
\end{center}

Exact substitution gives
\begin{equation}\label{eq:sextic}
 P_G(q)=\fall q{951531708}\bigl((q-479672508)^6-2\bigr).
\end{equation}
In particular, $\sqrt[6]{2}+479672508$ is a chromatic root.
The very large order is an artifact of this universal construction;
no claim of minimal order or minimal shift is made.

\subsection{A verification route independent of the partition sum}

Form the bipartite graph of cross-\emph{nonedges}, with its $d$ row
vertices on one side and $C$ on the other. Let $r_k$ be its number of
$k$-edge matchings. Coloring the two cliques with exactly $k$ reused
old colors gives
\begin{equation}\label{eq:rook}
 \frac{P_G(q)}{\fall qN}=\sum_{k=0}^d r_k\fall{q-N}{d-k}.
\end{equation}
Indeed, a matching chooses the rows using old colors and assigns
distinct old colors to them. The other $d-k$ rows choose distinct new
colors in $\fall{q-N}{d-k}$ ways. This is the matching formula used in
earlier biclique work \cite[Section 2.1]{Bohn}, and the counting argument
also proves it here.

For a type $J$ with $m_J$ vertices, assigning a chosen set $S\subseteq J$
of still-unmatched rows to that type can be done in
$\fall{m_J}{|S|}$ ways. A subset dynamic program over the row sets
therefore computes all $r_k$ directly from the atoms. The verifier
checks the independent identity
\begin{equation}\label{eq:rook-verify}
 \sum_{k=0}^d r_k\fall{x-L}{d-k}=f(x).
\end{equation}
It does not reuse the partition polynomial or the affine correction
matrices.

The included implementation was run on $x^d-2$ and on a mixed-sign
coefficient polynomial in every degree $1\le d\le9$ (18 inputs).
All exact factor checks, positive-atom checks, positive-shift checks,
and matching checks passed. Separate implementations checked the same
mechanism through degree nine and further mixed inputs through degree
eight, including coefficients of magnitude up to $10^9$.
These finite checks support transcription and implementation accuracy;
the universal assertion is proved by the preceding lemmas.

\section{Claim boundaries and literature status}\label{sec:status}

The original statement is kept unchanged: arbitrary algebraic integers,
positive integer translation, and zeros of ordinary unweighted chromatic
polynomials of finite simple graphs. The affine coordinates used here
are accompanied by the exact identity \eqref{eq:graph}. The coefficient
completion and the $t^2$ versus $\htf t$ estimate provide an actual
positive integral graph for every input.

As checked on 28 November 2026, Cameron's conjecture page, dated
2 November 2026, still lists the $\alpha+n$ conjecture among the
conjectures he would like to see settled \cite{CameronWeb}. The search
and primary sources consulted did not identify an earlier unrestricted
all-degree theorem or the same proof. This is evidence about the
consulted literature, not a proof of absolute priority.

The reviewed full texts were Cameron--Morgan (2017), Bohn (2012),
Gilmore (2015), and Delbourgo--Morgan (2022).
For the 2019 article, the official abstract and bibliographic record
were reviewed; its full text was not available and was not reviewed.
This access limit is material to an exhaustive comparison of methods.
The corpus snapshot was fixed at commit
\begin{center}
\small\texttt{9f8547bfd3b8f6560f20c1b63028dc87a07a2e33}.
\end{center}
Manuscript 245 supplies the methodological connection stated in Section 1.

\begin{center}
\small
\begin{tabularx}{\linewidth}{@{}l X@{}}
\toprule
Criterion & What this manuscript establishes or documents\\
\midrule
G1: named predecessor & Cameron--Morgan Conjecture 4 and Gilmore
Conjecture 7.1 are identified in primary texts.\\[3pt]
G2: unchanged target & Equation~\eqref{eq:main} implies the original
translation statement without changing $\alpha$ or restricting its degree.\\[3pt]
G3: prior open status & The current author-maintained problem list and
the inspected papers support the recorded status; exhaustive historical
absence is not certified.\\[3pt]
G4: unconditional proof & Sections 2--6 give the complete construction,
integer completion, positivity estimates, and termination argument for
all inputs, with no conjectural hypothesis.\\[3pt]
G5: substantive extension & The Hermite interior, all-degree basis rank,
factorial lattice correction, and uniform size control provide
unrestricted degree and input coefficients.\\[3pt]
G6: originality and priority & No identical resolution was identified
in the sources searched. Absence of an earlier identical theorem or
essentially identical proof is not independently certified.\\
\bottomrule
\end{tabularx}
\end{center}

The document is a research manuscript with internal proof audits and
exact computational checks. It has not received external peer review
or proof-assistant verification. Thus it records a complete proposed
mathematical proof while leaving independent validation, exhaustive
novelty assessment, and priority determination open.

\appendix
\section{An optional rational version of the leading profile}
\label{app:rational-profile}

This appendix is not needed for the main theorem. It shows that the
Hermite profile can, if desired, be replaced by a fixed strictly positive
rational profile for each degree. Algebraic-root isolation is then needed
only during a degree-dependent preliminary construction.

\begin{proposition}[A rational profile in every degree]
\label{prop:rational-profile}
Fix $d\ge2$, put $m=d-1$ and $\gamma=2^{2-d}$, and let
\[
 h(X)=\operatorname{He}_d(X)
     =\exp\!\left(-\frac12\frac{d^2}{dX^2}\right)X^d.
\]
Choose distinct rational numbers $u_1,\ldots,u_m$, set
$u_d=-\sum_{i=1}^m u_i$, and define
\begin{equation}\label{eq:rational-star}
 p(X)=\prod_{i=1}^m(X+u_i),\qquad
 v_i=1+\frac{h(-u_i)}{p'(-u_i)}.
\end{equation}
Give every outer pair $ij\subseteq[m]$ weight $1$ and the pair $id$
weight $v_i$. Then the weighted matching polynomial
\[
 \sum_{\mathcal M\text{ a matching of }K_d}
       \prod_{ij\in\mathcal M}v_{ij}
       \prod_{k\text{ unmatched}}(X+u_k)
\]
is exactly $X^d$.
Moreover, the rational $u_i$ can be chosen so that these pair weights
extend to strictly positive rational palette atoms whose singleton
moments are $2$. The outer pair weights remain $1$, and all leading
coefficient matrices required by the integral correction are nonsingular.
\end{proposition}

\begin{proof}
Write $T=\exp(\partial_X^2/2)$. A matching either leaves row $d$ unmatched
or pairs it with one of the outer rows. The matching polynomial is therefore
\[
 (X+u_d)Tp+\sum_{i=1}^m v_i\,T\!\left(\frac{p}{X+u_i}\right).
\]
The identity $T^{-1}XT=X-\partial_X$ shows that this equals $X^d$ if and
only if
\begin{equation}\label{eq:rational-interpolation}
 h=(X+u_d)p-p'+\sum_{i=1}^m v_i\frac{p}{X+u_i}.
\end{equation}
Both sides of \eqref{eq:rational-interpolation} are monic of degree $d$
and have zero coefficient of $X^{d-1}$. Their difference thus has degree
at most $d-2=m-1$. At $X=-u_i$, the right side has value
$(v_i-1)p'(-u_i)=h(-u_i)$ by \eqref{eq:rational-star}.
The difference vanishes at $m$ distinct points, proving the identity.
All the numbers $v_i$ are rational.

Let $\lambda_1,\ldots,\lambda_d$ be the distinct real roots of
$\operatorname{He}_d$. Choose the rational outer numbers $u_i$ close to
$-\lambda_i$, for $i\le m$. Their distinctness persists in a sufficiently
small neighborhood. Since the sum of the Hermite roots is zero,
$u_d$ approaches $-\lambda_d$. Formula \eqref{eq:rational-star} is
continuous there, with nonzero denominators, and every $v_i$ approaches
$1$. We may consequently arrange
\begin{equation}\label{eq:rational-smallness}
 \delta_i:=v_i-1,\qquad
 |\delta_i|<\frac{\gamma}{2m}\quad(1\le i\le m).
\end{equation}

Start with weight $\gamma$ on each of the $2^d-1$ nonempty atoms.
Its singleton moments are $2$, and its moments on a set $B$ of size
at least two are $2^{2-|B|}$; in particular its pair deficits are $1$
when the leading value of $L$ is $2$. Change only the star-pair deficits
from $1$ to $v_i$. Boolean inversion gives the following atom weights:
\[
 \begin{aligned}
 w_{\{i,d\}}&=\gamma-\delta_i,&
 w_{\{i\}}&=\gamma+\delta_i\qquad(i\le m),\\
 w_{\{d\}}&=\gamma+\sum_{i=1}^m\delta_i,&
 w_J&=\gamma\quad\text{for every other nonempty }J.
 \end{aligned}
\]
They are all at least $\gamma/2>0$ by
\eqref{eq:rational-smallness}. The changes in each singleton moment
cancel, so all those moments remain $2$. All outer data remain the
uniform data. The total atom mass is
$4-\gamma+\sum_i\delta_i>2$, as required for a positive leading shift.

Finally, for any $r+1$ distinct outer indices $J$, the polynomials
\[
 \prod_{j\in J\setminus\{i\}}(X+u_j),\qquad i\in J,
\]
form a basis of the polynomials of degree at most $r$: evaluation at
$-u_i$ isolates its $i$th member. Applying the invertible operator $T$
preserves their independence. These are exactly the leading complementary
palette polynomials used in the correction matrices, because every
outer pair weight remains $1$.
\end{proof}

The choice is effective. Successively refine rational isolating intervals
for the simple Hermite roots, choose distinct rational outer values, and
compute \eqref{eq:rational-star} exactly until
\eqref{eq:rational-smallness} holds. Continuity proves termination, and
the acceptance test uses rational arithmetic only. A common denominator
of this fixed profile can then be incorporated into the admissible scale.
This preliminary search depends on $d$, not on the coefficients of the
polynomial subsequently realized.


\section{Reference implementation}\label{app:code}

The following code is included literally in this single source file.
It was executed with Python 3.12.14 and SymPy 1.14.0. SymPy supplies
rational isolating intervals for the explicit auxiliary Hermite
polynomial; Gaussian elimination, lattice rounding, Boolean inversion,
and matching verification use Python integers and
\texttt{fractions.Fraction}. Copy the listing to a Python file and run
it to reproduce the complete sextic certificate. For another input,
call \texttt{construct([a1,...,ad])} followed by
\texttt{verify\_matching(result)}. The scale loop has no fixed cap.
The returned atoms give a compressed graph description, avoiding
expansion of its potentially enormous vertex set.

\begin{lstlisting}
"""Integral palette lifting in arbitrary degree.

Input [a1,...,ad] represents x**d+a1*x**(d-1)+...+ad.
SymPy supplies rational Hermite root intervals only. All construction,
lattice correction and certificate arithmetic is exact.
"""
from fractions import Fraction as Q
from functools import lru_cache
from math import comb, factorial
import sympy as sp


def multiply(p, q):
    out = [0] * (len(p) + len(q) - 1)
    for i, a in enumerate(p):
        for j, b in enumerate(q):
            out[i + j] += a * b
    return out


def add_scaled(p, q, c=1):
    out = list(p) + [0] * max(0, len(q) - len(p))
    for j, b in enumerate(q):
        out[j] += c * b
    return out


def solve(columns, rhs):
    n = len(rhs)
    a = [[Q(columns[j][i]) for j in range(n)] + [Q(rhs[i])]
         for i in range(n)]
    for j in range(n):
        pivot = next(i for i in range(j, n) if a[i][j])
        a[j], a[pivot] = a[pivot], a[j]
        v = a[j][j]
        a[j] = [x / v for x in a[j]]
        for i in range(n):
            if i != j:
                v = a[i][j]
                a[i] = [x - v * y for x, y in zip(a[i], a[j])]
    return [row[-1] for row in a]


@lru_cache(None)
def hermite(d):
    x = sp.Symbol("x")
    p, q = sp.Integer(1), x
    for j in range(1, d):
        p, q = q, sp.expand(x * q - j * p)
    return sp.Poly(q, x)


def singleton_values(d, t, M, a1):
    intervals = hermite(d).intervals(eps=sp.Rational(M, 4 * t))
    assert len(intervals) == d and all(e == 1 for _, e in intervals)
    z = []
    for i, ((lo, hi), _) in enumerate(intervals[:-1]):
        midpoint = (lo + hi) / 2
        u = Q(int(midpoint.p), int(midpoint.q))
        z.append(i + M * ((t * u - i) // M))
    z.append(a1 + comb(d, 2) - sum(z))
    return z


def construct_at(coeff, t):
    d = len(coeff)
    assert d >= 2 and all(isinstance(a, int) for a in coeff)
    M = factorial(d - 1)
    assert t > 0 and t % (2**d * M) == 0
    full = (1 << d) - 1
    center = 1 << (d - 1)
    outer = center - 1
    L = 2 * t * t
    z = singleton_values(d, t, M, coeff[0])
    deficits = {}
    for B in range(1, full + 1):
        k = B.bit_count()
        if k >= 2:
            assert t * t % (2**(k - 2)) == 0
            deficits[B] = L - t * t // (2**(k - 2))
            if not B & center:
                deficits[B] -= (B & -B).bit_length() - 1
    original = dict(deficits)

    @lru_cache(None)
    def outer_poly(S):
        assert not S & center
        if not S:
            return (1,)
        bit = S & -S
        tail = S ^ bit
        sub = tail
        out = [0] * (S.bit_count() + 1)
        while True:
            B = sub | bit
            k = B.bit_count()
            mu = (-1)**(k - 1) * factorial(k - 1)
            c = z[bit.bit_length() - 1] if k == 1 else -deficits[B]
            out = add_scaled(out, multiply([c, 1], outer_poly(S ^ B)), mu)
            if sub == 0:
                return tuple(out)
            sub = (sub - 1) & tail

    actual = multiply([z[-1], 1], outer_poly(outer))
    for B, D in deficits.items():
        if B & center:
            mu = (-1)**(B.bit_count() - 1) * factorial(B.bit_count() - 1)
            actual = add_scaled(actual, multiply([-D, 1], outer_poly(full ^ B)), mu)
    target = list(reversed(coeff)) + [1]

    def residual():
        return [a - b for a, b in zip(target, actual)]

    def derivative(B):
        k = B.bit_count()
        c = (-1)**k * factorial(k - 1)
        return [c * a for a in outer_poly(full ^ B)]

    def change(B, delta):
        nonlocal actual
        assert Q(delta).denominator == 1 and B & center
        delta = int(delta)
        deficits[B] += delta
        actual = add_scaled(actual, derivative(B), delta)

    assert residual()[d - 1:] == [0, 0]
    # The nested basis x^i (x+1)^(d-2-i) is unit triangular.
    r = residual()
    rho = []
    for j in range(d - 1):
        rho.append((r[j] - sum(comb(d - 2 - i, j - i) * rho[i]
                              for i in range(j))) % M)
    for i, v in enumerate(rho):
        change(center | (1 << i), v)
    assert all(a % M == 0 for a in residual())

    stage_maxima = []
    for k in range(2, d + 1):
        r = d - k
        I = (1 << (r + 1)) - 1
        complements = [I ^ (1 << i) for i in range(r + 1)]
        blocks = [full ^ R for R in complements]
        columns = [derivative(B) for B in blocks]
        error = residual()
        assert all(a == 0 for a in error[r + 1:])
        # Check the exact modular rank certificate at every layer.
        for i, R in enumerate(complements):
            for j, a in enumerate(outer_poly(R)):
                expected = comb(r - i, j - i) if j >= i else 0
                assert (a - expected) % M == 0
        exact = solve(columns, error[:r + 1])
        c = (-1)**k * factorial(k - 1)
        h = M // abs(c)
        rounded = [h * (v // h) for v in exact[:-1]]
        assert error[r] % c == 0
        rounded.append(error[r] // c - sum(rounded))
        assert all(v % h == 0 for v in rounded)
        for B, v in zip(blocks, rounded):
            change(B, v)
        assert all(a == 0 for a in residual()[r:])
        assert all(a % M == 0 for a in residual())
        stage_maxima.append(max(abs(v) for v in rounded))
    assert actual == target

    moments = [0] * (full + 1)
    for B in range(1, full + 1):
        moments[B] = L + z[B.bit_length() - 1] if B.bit_count() == 1 else L - deficits[B]
    atoms = list(moments)
    for i in range(d):
        for B in range(full + 1):
            if not B & (1 << i):
                atoms[B] -= atoms[B | (1 << i)]
    atoms = dict(enumerate(atoms[1:], 1))
    N = sum(atoms.values())
    return dict(degree=d, coefficients=list(coeff), t=t, M=M, L=L, z=z,
                deficits=deficits, atoms=atoms, empty_atom=0, N=N, shift=N-L,
                min_atom=min(atoms.values()), stage_maxima=stage_maxima,
                corrections={B: deficits[B]-original[B] for B in deficits
                             if deficits[B] != original[B]})


def construct(coeff):
    assert coeff and all(isinstance(a, int) for a in coeff)
    d = len(coeff)
    if d == 1:
        a = coeff[0]
        N = max(2, 2 - a)
        return dict(degree=1, coefficients=list(coeff), t=0, M=1, L=1-a,
                    z=[a], deficits={}, atoms={1: 1}, empty_atom=N-1,
                    N=N, shift=N+a-1, min_atom=1,
                    stage_maxima=[], corrections={})
    t = 2**d * factorial(d - 1)
    while True:
        result = construct_at(coeff, t)
        if result["min_atom"] > 0 and result["shift"] > 0:
            return result
        t *= 2


def verify_matching(w):
    """Count matchings in the nonedge bipartite graph independently."""
    d = w["degree"]
    full = (1 << d) - 1
    assert w["N"] == sum(w["atoms"].values()) + w["empty_atom"]
    assert w["shift"] > 0 and min(w["atoms"].values()) > 0
    dp = [0] * (full + 1)
    dp[0] = 1
    for J, count in w["atoms"].items():
        falling = [1]
        for k in range(1, d + 1):
            falling.append(falling[-1] * (count - k + 1))
        nxt = [0] * (full + 1)
        for used, ways in enumerate(dp):
            if not ways:
                continue
            available = J & (full ^ used)
            sub = available
            while True:
                nxt[used | sub] += ways * falling[sub.bit_count()]
                if sub == 0:
                    break
                sub = (sub - 1) & available
        dp = nxt
    rook = [sum(v for S, v in enumerate(dp) if S.bit_count() == k)
            for k in range(d + 1)]
    result = [0] * (d + 1)
    for k, number in enumerate(rook):
        basis = [number]
        for j in range(d - k):
            basis = multiply(basis, [-w["L"] - j, 1])
        result = add_scaled(result, basis)
    assert result == list(reversed(w["coefficients"])) + [1]
    return rook


if __name__ == "__main__":
    import json
    w = construct([0, 0, 0, 0, 0, -2])
    w["rook_numbers"] = verify_matching(w)
    print(json.dumps(w, indent=2))
\end{lstlisting}

\begin{thebibliography}{9}
\bibitem{CM}
P.~J. Cameron and K. Morgan,
\emph{Algebraic properties of chromatic roots},
Electronic Journal of Combinatorics \textbf{24}(1) (2017), P1.21.
\href{https://doi.org/10.37236/6578}{doi:10.37236/6578}.
\href{https://emis.de/ft/29131}{Full text}.

\bibitem{Bohn}
A. Bohn, \emph{Chromatic roots as algebraic integers},
Discrete Mathematics and Theoretical Computer Science Proceedings
\textbf{AR} (2012), 543--554.
\href{https://doi.org/10.46298/dmtcs.3061}{doi:10.46298/dmtcs.3061}.
\href{https://www.math.nagoya-u.ac.jp/fpsac12/download/contributed/dmAR0148.pdf}{Full text}.

\bibitem{Gilmore}
H.~J. Gilmore,
\emph{Algebraic Properties of Chromatic Polynomials and Their Roots},
MSc thesis, University of Waikato, 2015.
\url{https://hdl.handle.net/10289/9367}.

\bibitem{DM19}
D. Delbourgo and K. Morgan,
\emph{An algorithm which outputs a graph with a specified chromatic factor},
Discrete Applied Mathematics \textbf{257} (2019), 128--150.
\href{https://doi.org/10.1016/j.dam.2018.10.033}{doi:10.1016/j.dam.2018.10.033}.
Official abstract consulted through the author's
\href{https://profiles.auckland.ac.nz/daniel-delbourgo/publications}{University of Auckland publications record};
full text not reviewed.

\bibitem{DM22}
D. Delbourgo and K. Morgan, \emph{On the fifth chromatic coefficient},
Australasian Journal of Combinatorics \textbf{84}(1) (2022), 56--85.
\url{https://ajc.maths.uq.edu.au/pdf/84/ajc_v84_p056.pdf}.

\bibitem{VMRC}
Verifiable Mathematical Reasoning Corpus, manuscript 245,
\emph{Roots and coefficient peaks of graph polynomials},
research manuscript dated 21 November 2026, Sections 2--3.
Snapshot \texttt{9f8547bfd3b8f6560f20c1b63028dc87a07a2e33},
accessed 28 November 2026.
\url{https://github.com/QEEDDD/Verifiable-Mathematical-Reasoning-Corpus}.

\bibitem{CameronWeb}
P.~J. Cameron, \emph{Peter Cameron's conjectures},
page dated 2 November 2026, accessed 28 November 2026.
\url{https://cameroncounts.github.io/web/problems/cameronconjs.html}.
\end{thebibliography}
\end{document}
